\begin{recipe}{Focaccia}
\label{recipe:focaccia}
\kicker[Brood,Italiaans,Basisrecept]{Brood \& basis · Italiaans · basisrecept}
\index[register]{Focaccia}
\index[register]{Bloem}
\index[register]{Rozemarijn}
\index[register]{Gist}

\meta{Rijzen ca. 3 uur \dvd 1 grote focaccia \dvd Oven 220\,°C}

\lettrine{E}{en} nat, plakkerig deeg is hier geen probleem maar juist de
bedoeling: hoe natter het deeg, hoe luchtiger de focaccia. Het meeste werk is
wachten tot de gist zijn werk heeft gedaan.

\blockrule{Ingrediënten}
\begin{ingredients}
  \ing{500 g}{bloem type 00 (tipo 00)}\index[register]{Bloem}
  \ing{400 ml}{lauwwarm water (80\% hydratatie)}
  \ing{2 el}{extra vergine olijfolie (voor het deeg, ca.\ 5,5\%)}
  \ing{10 g}{zout (2\%)}
  \ing{7 g}{droge gist}
  \ing{3–4 el}{extra vergine olijfolie (voor bovenop en de bakvorm)}
  \ing{2 el}{water (voor de salamoia)}
  \ingb{grof zeezout}
  \ingb{rozemarijn}\index[register]{Rozemarijn}
\end{ingredients}

\blockrule{Bereiding}
\tip{Voor extra smaak kun je het deeg na de eerste rijs een nacht afgedekt in
de koelkast laten fermenteren. Haal het dan ongeveer een uur voor het bakken
uit de koelkast en ga verder met de bakvorm.}

\medskip
\noindent\textit{Deeg maken}
\begin{steps}
  \item Meng het lauwwarme water met de gist.
  \item Voeg de bloem toe en meng tot alles is opgenomen.
  \item Voeg het zout en 2 eetlepels olijfolie toe.
  \item Kneed 8–10 minuten tot een soepel, plakkerig deeg.
\end{steps}

\medskip
\noindent\textit{Eerste rijs}
\begin{steps}
  \item Doe het deeg in een ingevette kom, dek af en laat 1,5–2 uur rijzen,
        tot het ongeveer verdubbeld is.
\end{steps}

\medskip
\noindent\textit{In de bakvorm}
\begin{steps}
  \item Vet een bakblik flink in met olijfolie.
  \item Leg het deeg erin en rek het voorzichtig uit tot de randen van de
        vorm. Lukt dat nog niet? Laat het 15 minuten rusten en probeer
        opnieuw.
\end{steps}

\medskip
\noindent\textit{Tweede rijs}
\begin{steps}
  \item Laat het deeg in de vorm nog 30–45 minuten rijzen. Verwarm
        ondertussen de oven voor op 220\,°C (boven- en onderwarmte), zodat
        hij op temperatuur is als het deeg klaar is.
\end{steps}

\medskip
\noindent\textit{Afwerken en bakken}
\begin{steps}
  \item Besprenkel het deeg ruim met olijfolie en druk met natte of
        ingevette vingers kuiltjes in het deeg.
  \item Meng 2 eetlepels water met 2 eetlepels olijfolie en een snuf grof
        zeezout tot een salamoia en giet dat over de kuiltjes: het water
        verdampt in de oven, de olie frituurt de bovenkant licht en de
        kuiltjes blijven sappig, precies het verschil met "brood met olie
        erop".
  \item Bestrooi met grof zeezout en rozemarijn.
  \item Bak 20–25 minuten, tot de focaccia goudbruin is.
\end{steps}

\medskip
\noindent\textit{Serveren}
\begin{steps}
  \item Laat 10 minuten afkoelen op een rooster.
  \item Besprenkel eventueel nog met een klein scheutje olijfolie.
\end{steps}

\heroimagefade{images/focaccia}
\end{recipe}
